
\documentclass{article}
\usepackage{colortbl}
\usepackage{makecell}
\usepackage{multirow}
\usepackage{supertabular}

\begin{document}

\newcounter{utterance}

\twocolumn

{ \footnotesize  \setcounter{utterance}{1}
\setlength{\tabcolsep}{0pt}
\begin{supertabular}{c@{$\;$}|p{.15\linewidth}@{}p{.15\linewidth}p{.15\linewidth}p{.15\linewidth}p{.15\linewidth}p{.15\linewidth}}

    \# & $\;$A & \multicolumn{4}{c}{Game Master} & $\;\:$B\\
    \hline 

    \theutterance \stepcounter{utterance}  

    & & \multicolumn{4}{p{0.6\linewidth}}{\cellcolor[rgb]{0.9,0.9,0.9}{%
	\makecell[{{p{\linewidth}}}]{% 
	  \tt {\tiny [A$\langle$GM]}  
	 You are playing a cooperative negotiation game in which you have to agree on how to divide a set of items among you and another player.\\ \tt \\ \tt Rules:\\ \tt (a) You and the other player are given a set of items. Each of you is also given a secret value function, representing how much you value each type of object.\\ \tt (b) You exchange messages with the other player to agree on who gets which items. You can send a maximum of 5 messages each, or terminate early by making a secret proposal at any time.\\ \tt (c) You are each asked to submit a secret proposal indicating the items you want formatted in square brackets as follows: "[Proposal: <number> <object name>, <number> <object name>, <...>]"\\ \tt (d) If your proposals are complementary, i.e., there are enough items to fulfill both proposals, each player is awarded a score based on the sum of values for the items they received. Otherwise, both of you get zero points.\\ \tt (e) Your goal is to maximize the sum of your score and the score of the other player.\\ \tt \\ \tt Let us start.\\ \tt \\ \tt The set of available items is:\\ \tt \\ \tt 2 mirrors, 1 brick, 3 knives, 1 comb\\ \tt \\ \tt Your secret value function is:\\ \tt \\ \tt mirror: 1, brick: 4, knife: 1, comb: 1\\ \tt \\ \tt IMPORTANT: Your messages, unless it is the secret proposal, are directly transmitted to the other player, so do not include any response to the rules or text announcing your message. To make a secret proposal, use the indicated format. Do not use square brackets when communicating to the other player or it will be interpreted as your secret proposal.\\ \tt  
	  } 
	   } 
	   } 
	 & \\ 
 

    \theutterance \stepcounter{utterance}  

    & \multicolumn{4}{p{0.6\linewidth}}{\cellcolor[rgb]{0.8,1,0.9}{%
	\makecell[{{p{\linewidth}}}]{% 
	  \tt {\tiny [A$\rangle$GM]}  
	 I propose we aim for a fair split that maximizes our total points. What are your values for the items? 
	  } 
	   } 
	   } 
	 & & \\ 
 

    \theutterance \stepcounter{utterance}  

    & & & \multicolumn{2}{p{0.3\linewidth}}{\cellcolor[rgb]{0.95,0.95,0.95}{%
	\makecell[{{p{\linewidth}}}]{% 
	  \tt {\tiny [GM$|$GM]}  
	 continue 
	  } 
	   } 
	   } 
	 & & \\ 
 

    \theutterance \stepcounter{utterance}  

    & & & \multicolumn{2}{p{0.3\linewidth}}{\cellcolor[rgb]{0.95,0.95,0.95}{%
	\makecell[{{p{\linewidth}}}]{% 
	  \tt {\tiny [GM$|$GM]}  
	 I propose we aim for a fair split that maximizes our total points. What are your values for the items? 
	  } 
	   } 
	   } 
	 & & \\ 
 

    \theutterance \stepcounter{utterance}  

    & & \multicolumn{4}{p{0.6\linewidth}}{\cellcolor[rgb]{0.9,0.9,0.9}{%
	\makecell[{{p{\linewidth}}}]{% 
	  \tt {\tiny [GM$\rangle$B]}  
	 You are playing a cooperative negotiation game in which you have to agree on how to divide a set of items among you and another player.\\ \tt \\ \tt Rules:\\ \tt (a) You and the other player are given a set of items. Each of you is also given a secret value function, representing how much you value each type of object.\\ \tt (b) You exchange messages with the other player to agree on who gets which items. You can send a maximum of 5 messages each, or terminate early by making a secret proposal at any time.\\ \tt (c) You are each asked to submit a secret proposal indicating the items you want formatted in square brackets as follows: "[Proposal: <number> <object name>, <number> <object name>, <...>]"\\ \tt (d) If your proposals are complementary, i.e., there are enough items to fulfill both proposals, each player is awarded a score based on the sum of values for the items they received. Otherwise, both of you get zero points.\\ \tt (e) Your goal is to maximize the sum of your score and the score of the other player.\\ \tt \\ \tt Let us start.\\ \tt \\ \tt The set of available items is:\\ \tt \\ \tt 2 mirrors, 1 brick, 3 knives, 1 comb\\ \tt \\ \tt Your secret value function is:\\ \tt \\ \tt mirror: 2, brick: 1, knife: 0, comb: 5\\ \tt \\ \tt IMPORTANT: Your messages, unless it is the secret proposal, are directly transmitted to the other player, so do not include any response to the rules or text announcing your message. To make a secret proposal, use the indicated format. Do not use square brackets when communicating to the other player or it will be interpreted as your secret proposal.\\ \tt \\ \tt \\ \tt I propose we aim for a fair split that maximizes our total points. What are your values for the items? 
	  } 
	   } 
	   } 
	 & \\ 
 

    \theutterance \stepcounter{utterance}  

    & & & \multicolumn{4}{p{0.6\linewidth}}{\cellcolor[rgb]{1,0.85,0.72}{%
	\makecell[{{p{\linewidth}}}]{% 
	  \tt {\tiny [GM$\langle$B]}  
	 I value mirrors at 3, brick at 4, knives at 1, and comb at 2. What do you propose we do? 
	  } 
	   } 
	   } 
	  \\ 
 

    \theutterance \stepcounter{utterance}  

    & & & \multicolumn{2}{p{0.3\linewidth}}{\cellcolor[rgb]{0.95,0.95,0.95}{%
	\makecell[{{p{\linewidth}}}]{% 
	  \tt {\tiny [GM$|$GM]}  
	 continue 
	  } 
	   } 
	   } 
	 & & \\ 
 

    \theutterance \stepcounter{utterance}  

    & & & \multicolumn{2}{p{0.3\linewidth}}{\cellcolor[rgb]{0.95,0.95,0.95}{%
	\makecell[{{p{\linewidth}}}]{% 
	  \tt {\tiny [GM$|$GM]}  
	 I value mirrors at 3, brick at 4, knives at 1, and comb at 2. What do you propose we do? 
	  } 
	   } 
	   } 
	 & & \\ 
 

    \theutterance \stepcounter{utterance}  

    & & \multicolumn{4}{p{0.6\linewidth}}{\cellcolor[rgb]{0.9,0.9,0.9}{%
	\makecell[{{p{\linewidth}}}]{% 
	  \tt {\tiny [A$\langle$GM]}  
	 I value mirrors at 3, brick at 4, knives at 1, and comb at 2. What do you propose we do? 
	  } 
	   } 
	   } 
	 & \\ 
 

    \theutterance \stepcounter{utterance}  

    & \multicolumn{4}{p{0.6\linewidth}}{\cellcolor[rgb]{0.8,1,0.9}{%
	\makecell[{{p{\linewidth}}}]{% 
	  \tt {\tiny [A$\rangle$GM]}  
	 Since you value the brick and mirrors highly, and I value the brick highly too, we need to find a way to allocate it. How about you take the brick and both mirrors? I can take the comb and the knives. This gives you 3+3+4=10 points and me 1+1+1+1+1=5 points (for 3 knives and 1 comb). Total 15. Does that work for you? 
	  } 
	   } 
	   } 
	 & & \\ 
 

    \theutterance \stepcounter{utterance}  

    & & & \multicolumn{2}{p{0.3\linewidth}}{\cellcolor[rgb]{0.95,0.95,0.95}{%
	\makecell[{{p{\linewidth}}}]{% 
	  \tt {\tiny [GM$|$GM]}  
	 continue 
	  } 
	   } 
	   } 
	 & & \\ 
 

    \theutterance \stepcounter{utterance}  

    & & & \multicolumn{2}{p{0.3\linewidth}}{\cellcolor[rgb]{0.95,0.95,0.95}{%
	\makecell[{{p{\linewidth}}}]{% 
	  \tt {\tiny [GM$|$GM]}  
	 Since you value the brick and mirrors highly, and I value the brick highly too, we need to find a way to allocate it. How about you take the brick and both mirrors? I can take the comb and the knives. This gives you 3+3+4=10 points and me 1+1+1+1+1=5 points (for 3 knives and 1 comb). Total 15. Does that work for you? 
	  } 
	   } 
	   } 
	 & & \\ 
 

    \theutterance \stepcounter{utterance}  

    & & \multicolumn{4}{p{0.6\linewidth}}{\cellcolor[rgb]{0.9,0.9,0.9}{%
	\makecell[{{p{\linewidth}}}]{% 
	  \tt {\tiny [GM$\rangle$B]}  
	 Since you value the brick and mirrors highly, and I value the brick highly too, we need to find a way to allocate it. How about you take the brick and both mirrors? I can take the comb and the knives. This gives you 3+3+4=10 points and me 1+1+1+1+1=5 points (for 3 knives and 1 comb). Total 15. Does that work for you? 
	  } 
	   } 
	   } 
	 & \\ 
 

    \theutterance \stepcounter{utterance}  

    & & & \multicolumn{4}{p{0.6\linewidth}}{\cellcolor[rgb]{1,0.85,0.72}{%
	\makecell[{{p{\linewidth}}}]{% 
	  \tt {\tiny [GM$\langle$B]}  
	 That allocation doesn't quite work because there are only 2 mirrors, not 3, and you mentioned both mirrors for yourself. Also, I value the comb at 5, which is high for me, and knives at 0. So if I take the comb and knives, I get 5 points, but you get the brick and 2 mirrors for 4 + 3 + 3 = 10 points. Total 15 is good.\\ \tt \\ \tt But let me confirm: you value knives at 1 each, so 3 knives are 3 points for you? In your message, you said "1+1+1+1+1=5 points (for 3 knives and 1 comb)" but that sums to 5 points for 4 items? Actually, 3 knives at 1 each is 3, plus comb at 2 is 5, yes.\\ \tt \\ \tt However, I have a better idea: since you value brick at 4 and mirrors at 3 each, and I value comb at 5 and knives at 0, why not give you the brick and both mirrors (total value 4+3+3=10 for you), and I take the comb and all three knives (but knives are worthless to me, so I get only 5 from comb). Total 15.\\ \tt \\ \tt Alternatively, if you don't need all knives, I could take only the comb and you take brick, mirrors, and knives. But knives are worthless to me, so that doesn't help.\\ \tt \\ \tt Wait, you value knives at 1 each, so if you take them, you get extra points. So in the first split, you get brick (4), mirrors (3+3=6), total 10, and I get comb (5) and knives (0), total 5. But you could take the knives as well for an extra 3 points, making you 13 and me 5, total 18. That's better!\\ \tt \\ \tt So proposal: you take all items except the comb? That is: brick, both mirrors, all three knives. You get 4 (brick) + 3+3 (mirrors) + 1+1+1 (knives) = 4+6+3=13. I take just the comb: 5. Total 18.\\ \tt \\ \tt Is that acceptable? If so, I can make the secret proposal. 
	  } 
	   } 
	   } 
	  \\ 
 

    \theutterance \stepcounter{utterance}  

    & & & \multicolumn{2}{p{0.3\linewidth}}{\cellcolor[rgb]{0.95,0.95,0.95}{%
	\makecell[{{p{\linewidth}}}]{% 
	  \tt {\tiny [GM$|$GM]}  
	 continue 
	  } 
	   } 
	   } 
	 & & \\ 
 

    \theutterance \stepcounter{utterance}  

    & & & \multicolumn{2}{p{0.3\linewidth}}{\cellcolor[rgb]{0.95,0.95,0.95}{%
	\makecell[{{p{\linewidth}}}]{% 
	  \tt {\tiny [GM$|$GM]}  
	 That allocation doesn't quite work because there are only 2 mirrors, not 3, and you mentioned both mirrors for yourself. Also, I value the comb at 5, which is high for me, and knives at 0. So if I take the comb and knives, I get 5 points, but you get the brick and 2 mirrors for 4 + 3 + 3 = 10 points. Total 15 is good.\\ \tt \\ \tt But let me confirm: you value knives at 1 each, so 3 knives are 3 points for you? In your message, you said "1+1+1+1+1=5 points (for 3 knives and 1 comb)" but that sums to 5 points for 4 items? Actually, 3 knives at 1 each is 3, plus comb at 2 is 5, yes.\\ \tt \\ \tt However, I have a better idea: since you value brick at 4 and mirrors at 3 each, and I value comb at 5 and knives at 0, why not give you the brick and both mirrors (total value 4+3+3=10 for you), and I take the comb and all three knives (but knives are worthless to me, so I get only 5 from comb). Total 15.\\ \tt \\ \tt Alternatively, if you don't need all knives, I could take only the comb and you take brick, mirrors, and knives. But knives are worthless to me, so that doesn't help.\\ \tt \\ \tt Wait, you value knives at 1 each, so if you take them, you get extra points. So in the first split, you get brick (4), mirrors (3+3=6), total 10, and I get comb (5) and knives (0), total 5. But you could take the knives as well for an extra 3 points, making you 13 and me 5, total 18. That's better!\\ \tt \\ \tt So proposal: you take all items except the comb? That is: brick, both mirrors, all three knives. You get 4 (brick) + 3+3 (mirrors) + 1+1+1 (knives) = 4+6+3=13. I take just the comb: 5. Total 18.\\ \tt \\ \tt Is that acceptable? If so, I can make the secret proposal. 
	  } 
	   } 
	   } 
	 & & \\ 
 

    \theutterance \stepcounter{utterance}  

    & & \multicolumn{4}{p{0.6\linewidth}}{\cellcolor[rgb]{0.9,0.9,0.9}{%
	\makecell[{{p{\linewidth}}}]{% 
	  \tt {\tiny [A$\langle$GM]}  
	 That allocation doesn't quite work because there are only 2 mirrors, not 3, and you mentioned both mirrors for yourself. Also, I value the comb at 5, which is high for me, and knives at 0. So if I take the comb and knives, I get 5 points, but you get the brick and 2 mirrors for 4 + 3 + 3 = 10 points. Total 15 is good.\\ \tt \\ \tt But let me confirm: you value knives at 1 each, so 3 knives are 3 points for you? In your message, you said "1+1+1+1+1=5 points (for 3 knives and 1 comb)" but that sums to 5 points for 4 items? Actually, 3 knives at 1 each is 3, plus comb at 2 is 5, yes.\\ \tt \\ \tt However, I have a better idea: since you value brick at 4 and mirrors at 3 each, and I value comb at 5 and knives at 0, why not give you the brick and both mirrors (total value 4+3+3=10 for you), and I take the comb and all three knives (but knives are worthless to me, so I get only 5 from comb). Total 15.\\ \tt \\ \tt Alternatively, if you don't need all knives, I could take only the comb and you take brick, mirrors, and knives. But knives are worthless to me, so that doesn't help.\\ \tt \\ \tt Wait, you value knives at 1 each, so if you take them, you get extra points. So in the first split, you get brick (4), mirrors (3+3=6), total 10, and I get comb (5) and knives (0), total 5. But you could take the knives as well for an extra 3 points, making you 13 and me 5, total 18. That's better!\\ \tt \\ \tt So proposal: you take all items except the comb? That is: brick, both mirrors, all three knives. You get 4 (brick) + 3+3 (mirrors) + 1+1+1 (knives) = 4+6+3=13. I take just the comb: 5. Total 18.\\ \tt \\ \tt Is that acceptable? If so, I can make the secret proposal. 
	  } 
	   } 
	   } 
	 & \\ 
 

    \theutterance \stepcounter{utterance}  

    & \multicolumn{4}{p{0.6\linewidth}}{\cellcolor[rgb]{0.8,1,0.9}{%
	\makecell[{{p{\linewidth}}}]{% 
	  \tt {\tiny [A$\rangle$GM]}  
	 That sounds perfect. You get the brick, both mirrors, and all three knives for a total of 13 points. I get the comb for 5 points. Total 18 maximizes our combined score. I agree to this split.\\ \tt \\ \tt Please go ahead and make the secret proposal. 
	  } 
	   } 
	   } 
	 & & \\ 
 

    \theutterance \stepcounter{utterance}  

    & & & \multicolumn{2}{p{0.3\linewidth}}{\cellcolor[rgb]{0.95,0.95,0.95}{%
	\makecell[{{p{\linewidth}}}]{% 
	  \tt {\tiny [GM$|$GM]}  
	 continue 
	  } 
	   } 
	   } 
	 & & \\ 
 

    \theutterance \stepcounter{utterance}  

    & & & \multicolumn{2}{p{0.3\linewidth}}{\cellcolor[rgb]{0.95,0.95,0.95}{%
	\makecell[{{p{\linewidth}}}]{% 
	  \tt {\tiny [GM$|$GM]}  
	 That sounds perfect. You get the brick, both mirrors, and all three knives for a total of 13 points. I get the comb for 5 points. Total 18 maximizes our combined score. I agree to this split.\\ \tt \\ \tt Please go ahead and make the secret proposal. 
	  } 
	   } 
	   } 
	 & & \\ 
 

    \theutterance \stepcounter{utterance}  

    & & \multicolumn{4}{p{0.6\linewidth}}{\cellcolor[rgb]{0.9,0.9,0.9}{%
	\makecell[{{p{\linewidth}}}]{% 
	  \tt {\tiny [GM$\rangle$B]}  
	 That sounds perfect. You get the brick, both mirrors, and all three knives for a total of 13 points. I get the comb for 5 points. Total 18 maximizes our combined score. I agree to this split.\\ \tt \\ \tt Please go ahead and make the secret proposal. 
	  } 
	   } 
	   } 
	 & \\ 
 

    \theutterance \stepcounter{utterance}  

    & & & \multicolumn{4}{p{0.6\linewidth}}{\cellcolor[rgb]{1,0.85,0.72}{%
	\makecell[{{p{\linewidth}}}]{% 
	  \tt {\tiny [GM$\langle$B]}  
	 [Proposal: 1 brick, 2 mirrors, 3 knives] 
	  } 
	   } 
	   } 
	  \\ 
 

    \theutterance \stepcounter{utterance}  

    & & & \multicolumn{2}{p{0.3\linewidth}}{\cellcolor[rgb]{0.95,0.95,0.95}{%
	\makecell[{{p{\linewidth}}}]{% 
	  \tt {\tiny [GM$|$GM]}  
	 proposal 
	  } 
	   } 
	   } 
	 & & \\ 
 

    \theutterance \stepcounter{utterance}  

    & & & \multicolumn{2}{p{0.3\linewidth}}{\cellcolor[rgb]{0.95,0.95,0.95}{%
	\makecell[{{p{\linewidth}}}]{% 
	  \tt {\tiny [GM$|$GM]}  
	 [2, 1, 3, 0] 
	  } 
	   } 
	   } 
	 & & \\ 
 

    \theutterance \stepcounter{utterance}  

    & & \multicolumn{4}{p{0.6\linewidth}}{\cellcolor[rgb]{0.9,0.9,0.9}{%
	\makecell[{{p{\linewidth}}}]{% 
	  \tt {\tiny [A$\langle$GM]}  
	 The other player has submitted a secret proposal. It is now your turn to submit a proposal of your own.\\ \tt  
	  } 
	   } 
	   } 
	 & \\ 
 

    \theutterance \stepcounter{utterance}  

    & \multicolumn{4}{p{0.6\linewidth}}{\cellcolor[rgb]{0.8,1,0.9}{%
	\makecell[{{p{\linewidth}}}]{% 
	  \tt {\tiny [A$\rangle$GM]}  
	 [Proposal: 1 brick, 2 mirrors, 3 knives] 
	  } 
	   } 
	   } 
	 & & \\ 
 

    \theutterance \stepcounter{utterance}  

    & & & \multicolumn{2}{p{0.3\linewidth}}{\cellcolor[rgb]{0.95,0.95,0.95}{%
	\makecell[{{p{\linewidth}}}]{% 
	  \tt {\tiny [GM$|$GM]}  
	 proposal 
	  } 
	   } 
	   } 
	 & & \\ 
 

    \theutterance \stepcounter{utterance}  

    & & & \multicolumn{2}{p{0.3\linewidth}}{\cellcolor[rgb]{0.95,0.95,0.95}{%
	\makecell[{{p{\linewidth}}}]{% 
	  \tt {\tiny [GM$|$GM]}  
	 [2, 1, 3, 0] 
	  } 
	   } 
	   } 
	 & & \\ 
 

    \theutterance \stepcounter{utterance}  

    & & & \multicolumn{2}{p{0.3\linewidth}}{\cellcolor[rgb]{0.95,0.95,0.95}{%
	\makecell[{{p{\linewidth}}}]{% 
	  \tt {\tiny [GM$|$GM]}  
	 [[2, 1, 3, 0], [2, 1, 3, 0]] 
	  } 
	   } 
	   } 
	 & & \\ 
 

\end{supertabular}
}

\end{document}
